\documentclass[conference]{IEEEtran}
\IEEEoverridecommandlockouts

\usepackage{cite}
\usepackage{amsmath,amssymb,amsfonts}
\usepackage{graphicx}
\usepackage{textcomp}
\usepackage{xcolor}
\usepackage{booktabs}
\usepackage{multirow}
\usepackage{microtype}

\def\BibTeX{{\rm B\kern-.05em{\sc i\kern-.025em b}\kern-.08em
    T\kern-.1667em\lower.7ex\hbox{E}\kern-.125emX}}

\begin{document}

\title{MinImageScorer: Error-Driven Difficulty Scoring\\
for Guided Copy-Paste Augmentation in\\
Agricultural Object Detection}

% TODO: Fill in author/supervisor/affiliation
\author{
\IEEEauthorblockN{Author Name\textsuperscript{1},
                  Supervisor Name\textsuperscript{1}}
\IEEEauthorblockA{%
\textsuperscript{1}Department of XXX, University of XXX, Country\\
email@university.edu}
}

\maketitle

% ----------------------------------------------------------
\begin{abstract}
Copy-paste augmentation is widely applied to improve segmentation result
but it  
but is rarely evaluated against the specific failure mode a dataset
exhibits. We propose \textbf{MinImageScorer}, an error-driven framework
that scores per-object detection difficulty from baseline model failures
using a cost function combining prediction confidence and spatial
accuracy. We validate the score on two agricultural
benchmarks---MinneApple, Global Wheat Head (GWHD)---and use it to
guide copy-paste augmentation. Experiments across three random seeds
show that score-guided selection outperforms random selection on
MinneApple ($\Delta\mathrm{AP}{=}{+}0.013$) but neither condition
recovers the no-augmentation baseline. Feature-space analysis confirms
that augmented images shift the training distribution away from the
test distribution, explaining the consistent degradation. On GWHD,
domain shift dominates and copy-paste provides no benefit under either
selection strategy. Our findings establish that the utility of
copy-paste augmentation depends on the dataset's underlying failure
mode, and that MinImageScorer correctly identifies the objects
responsible for that failure.
\end{abstract}

\begin{IEEEkeywords}
object detection, copy-paste augmentation, difficulty scoring,
agricultural detection, small objects, domain shift
\end{IEEEkeywords}

% ----------------------------------------------------------
\section{Introduction}

Agricultural object detection benchmarks such as
MinneApple~\cite{hani2020minneapple} and Global Wheat Head
(GWHD)~\cite{david2021global} present distinct structural failure
modes. MinneApple is dominated by high-density small objects near the
resolution floor of standard backbones, while GWHD contains images
from multiple geographic sources with significant cross-domain visual
variation. These failure modes differ fundamentally from the object
\textit{underrepresentation} problem that Kisantal
\textit{et~al.}~\cite{kisantal2019augmentation} addressed on MS COCO,
where copy-paste augmentation produced measurable gains by increasing
small-object anchor matches during training. Ghiasi
\textit{et~al.}~\cite{ghiasi2021simple} further established
large-scale copy-paste as a general augmentation strategy for instance
segmentation, but likewise evaluated it on datasets without structural
resolution or domain-shift limitations.

A critical gap remains: practitioners routinely apply copy-paste
augmentation without first diagnosing whether the dataset's failure
mode is addressable by adding more object instances. When the failure
is structural---resolution-limited detectability or cross-domain
shift---copy-paste may corrupt the training distribution rather than
improve it. No prior work provides an empirical framework to diagnose
this mismatch before augmentation is applied.

We make the following contributions:
\begin{enumerate}
    \item \textbf{MinImageScorer}: a per-object difficulty score
    derived from model failure signals, with weights auto-calibrated
    via Pearson correlation optimization, requiring no domain-specific
    assumptions.
    \item \textbf{Empirical validation} that the score captures
    dataset-specific failure modes on both MinneApple and GWHD.
    \item \textbf{Controlled augmentation analysis} comparing
    score-guided and random copy-paste across three seeds, with
    feature-space analysis explaining the observed performance
    patterns.
\end{enumerate}

% ----------------------------------------------------------
\section{Methodology}

\subsection{MinImageScorer: Formulation}

Let $G(I)$ denote the ground-truth objects in image $I$ and
$\mathcal{P}$ the set of model predictions. For object $g$, let
$\mathcal{P}(g;\tau)$ be the predictions satisfying
$\mathrm{IoU}(p,g)\ge\tau$. The \emph{object difficulty score} is
defined as:

\begin{equation}
S_{\mathrm{obj}}(g) =
\begin{cases}
\displaystyle\min_{p\,\in\,\mathcal{P}(g;\tau)}
  \bigl[
    \alpha\,(1 - \mathrm{conf}(p))
    + \beta\,(1 - \mathrm{IoU}(p,g))
  \bigr],
  & \text{if } \mathcal{P}(g;\tau) \neq \varnothing, \\[6pt]
\alpha + \beta,
  & \text{otherwise.}
\end{cases}
\label{eq:obj_score}
\end{equation}

Missed detections receive the maximum penalty $\alpha + \beta$.
The \emph{image-level score} aggregates object difficulties and
penalizes false positives:

\begin{equation}
S_{\mathrm{img}}(I) =
  w_1 \cdot \frac{1}{|G(I)|} \sum_{g \in G(I)} S_{\mathrm{obj}}(g)
  \;+\; w_2 \cdot \mathrm{FPrate}(I;\,\gamma,\tau).
\label{eq:img_score}
\end{equation}

Weights $\alpha,\beta,w_1,w_2$ are calibrated by maximizing the
Pearson correlation between score distributions and per-image miss
rate and false-positive rate on the training set. This data-driven
calibration requires no domain-specific tuning.

\subsection{Score-Guided Copy-Paste Pipeline}

The pipeline proceeds in five steps:
(1)~train a baseline YOLOv8s~\cite{yolov8} on the original training
set with all built-in augmentations disabled;
(2)~run inference on the training set to collect per-prediction
confidence and IoU signals;
(3)~compute $S_{\mathrm{obj}}$ and $S_{\mathrm{img}}$ for all objects
and images;
(4)~sample objects proportionally to their scores using exponential
weighting ($\alpha{=}2$) for score-guided, or uniformly for random
selection; and
(5)~paste 2--8 objects per augmented image within the source image
(same-image pasting), then retrain from COCO pretrained weights on
the augmented set.
Both conditions are identical in all parameters except the selection
policy, ensuring a controlled comparison.

% ----------------------------------------------------------
\section{Experiments}

\subsection{Setup}

\textbf{Datasets.}
MinneApple~\cite{hani2020minneapple} contains
% ============================================================
% [FILL] Insert actual train/val/test image counts and total
%        object instance counts. Example format:
%        670 training, 67 validation, and 67 test images,
%        with 41,204 annotated apple instances.
% ============================================================
[FILL: N\textsubscript{train}/N\textsubscript{val}/N\textsubscript{test}]
images of apple trees with dense small-object annotations.
GWHD~\cite{david2021global} contains
% ============================================================
% [FILL] Insert GWHD train/val/test counts and domain count.
% ============================================================
[FILL: N\textsubscript{train}/N\textsubscript{val}/N\textsubscript{test}]
images of wheat fields collected across multiple geographic domains
(Australia, Japan, India, and others), making it a multi-domain
benchmark.

\textbf{Model and training.}
All models use YOLOv8s~\cite{yolov8} initialized from COCO pretrained
weights. Training settings: MinneApple 200 epochs / patience 10;
GWHD 100 epochs / patience 8; batch size 16; input size
$1024{\times}1024$; learning rate $10^{-2}$. All Ultralytics
built-in augmentations (mosaic, mixup, HSV jitter, flips) are
disabled to isolate the effect of copy-paste.

\textbf{Augmentation parameters.}
Dataset augmentation ratio 0.3; scale jitter $[0.9,\,1.1]$
(MinneApple) and $[0.8,\,1.2]$ (GWHD); no rotation; overlap
avoidance with 5\,px margin; no alpha blending.
Results are reported as means over seeds $\{123,456,789\}$.

\subsection{Results}

Table~\ref{tab:results} reports 3-seed mean COCO AP on original test
sets for both datasets.

\begin{table}[t]
\centering
\caption{Three-seed mean COCO AP on original test sets.
  \textbf{Bold} = best copy-paste condition per dataset.}
\label{tab:results}
\resizebox{\columnwidth}{!}{%
\begin{tabular}{@{}llrrrrr@{}}
\toprule
\textbf{Dataset} & \textbf{Condition}
  & \textbf{AP}
  & \textbf{AP\textsubscript{50}}
  & \textbf{AP\textsubscript{75}}
  & \textbf{AP\textsubscript{S}}
  & \textbf{AP\textsubscript{M}} \\
\midrule
\multirow{3}{*}{MinneApple}
  & Baseline          & 0.439 & 0.771 & 0.452 & 0.300 & 0.592 \\
  & Random CP         & 0.415 & 0.741 & 0.427 & 0.279 & 0.563 \\
  & \textbf{Score CP} & \textbf{0.428} & \textbf{0.759}
                      & \textbf{0.440} & \textbf{0.286} & \textbf{0.581} \\
\midrule
\multirow{3}{*}{GWHD}
  & Baseline          & 0.507 & 0.922 & 0.491 & 0.136 & 0.501 \\
  & Random CP         & 0.486 & 0.905 & 0.456 & 0.118 & 0.480 \\
  & \textbf{Score CP} & \textbf{0.484} & \textbf{0.905}
                      & \textbf{0.455} & \textbf{0.098} & \textbf{0.477} \\
\bottomrule
\end{tabular}%
}
\end{table}

On MinneApple, score-guided copy-paste outperforms random copy-paste
by $\Delta\mathrm{AP}{=}{+}0.013$, with consistent gains across
AP, AP$_{50}$, AP$_{75}$, AP$_S$, and AP$_M$. This indicates
the score successfully directs augmentation toward harder, smaller
objects that provide greater gradient signal. However, both copy-paste
conditions remain below baseline ($-0.024$ for Score CP), showing
that the distributional cost of pasting outweighs the selective
benefit on this dataset.

On GWHD, score-guided and random selection are
statistically indistinguishable (gap: $-0.002$~AP), and both degrade
below baseline ($-0.023$). Notably, on AP$_S$ the score-guided
condition is worse than random ($0.098$ vs.\ $0.118$), because
score-guided sampling concentrates pasted objects from
domain-minority sources (predominantly \textit{arvalis}), amplifying
the existing domain imbalance rather than correcting it.

\subsection{Feature-Space Analysis}

\begin{figure}[t]
\centering
% ============================================================
% [FILL] Replace the placeholder below with your actual figure.
% Recommended: export UMAP figure as feature_space_combined.pdf
% or .png, then uncomment and update the includegraphics line.
% ============================================================
% \includegraphics[width=\columnwidth]{feature_space_combined.pdf}
\framebox{\parbox{0.9\columnwidth}{\centering\vspace{1.8cm}
\small FIGURE PLACEHOLDER:\\
Image-level UMAP projections for MinneApple (left)\\
and GWHD (right). Groups: original train, augmented\\
(random CP), augmented (score CP), and test images.
\vspace{1.8cm}}}
\caption{Image-level backbone feature projections (YOLOv8s P4 layer,
  global average pool, UMAP). Augmented images occupy regions
  displaced from the test distribution relative to original training
  images on both datasets. Score-guided augmented images show greater
  displacement on GWHD due to domain-minority object concentration.}
\label{fig:feature_space}
\end{figure}

Fig.~\ref{fig:feature_space} shows image-level backbone embeddings
(YOLOv8s P4, global average pool, UMAP projection) for four groups:
original training images, random-CP augmented images, score-CP
augmented images, and test images.

On both datasets, augmented images occupy a region of feature space
displaced from the test distribution relative to the original training
images. Quantitatively, the centroid distance from augmented images to
the test set exceeds that of the original training images:
% ============================================================
% [FILL] Replace placeholders with actual centroid distances.
% Example: "MinneApple: 1.83 (Score CP) vs. 1.21 (original train);
%           GWHD: 2.14 (Score CP) vs. 1.56 (original train)."
% ============================================================
MinneApple: \textit{[value]} (Score~CP) vs.\ \textit{[value]}
(original train);
GWHD: \textit{[value]} (Score~CP) vs.\ \textit{[value]}
(original train).
On GWHD, score-guided augmented images show greater displacement than
random augmented images, consistent with domain-minority concentration.
This distributional shift directly explains the consistent performance
degradation observed in Table~\ref{tab:results}.

% ----------------------------------------------------------
\section{Discussion}

\textbf{Why copy-paste fails on these datasets.}
Kisantal~\textit{et~al.}~\cite{kisantal2019augmentation} demonstrated
that copy-paste improves AP$_S$ on MS COCO by increasing anchor
coverage for underrepresented small objects. This benefit presupposes
that (1) hard objects are above the backbone feature resolution floor,
and (2) pasted objects share the same visual domain as the target
scene. Neither condition holds reliably here. On MinneApple, the
hardest objects identified by MinImageScorer are predominantly
at or below the resolution limit of the backbone, making them weakly
represented in feature maps regardless of paste frequency. On GWHD,
even within-image pasting fails because the feature-space mismatch
originates at the image level. These structural limitations were not
present in the COCO setting studied by either Kisantal
\textit{et~al.}~\cite{kisantal2019augmentation} or Ghiasi
\textit{et~al.}~\cite{ghiasi2021simple}.

\textbf{Diagnostic value of the score.}
Despite producing no absolute improvement, MinImageScorer demonstrates
clear diagnostic utility: the score correlates with per-image miss
rate
% ============================================================
% [FILL] Insert Pearson r values from your calibration analysis.
% Example: "r = 0.71 on MinneApple, r = 0.58 on GWHD"
% ============================================================
($r{=}[\text{value}]$ on MinneApple, $r{=}[\text{value}]$ on GWHD),
and on MinneApple it selects objects that yield a $+0.013$~AP
relative gain over uniform random selection. The score identifies the
right objects; the limitation is that copy-paste cannot make those
objects more learnable when the failure mode is structural.

\textbf{Boundary condition for practitioners.}
Score-guided copy-paste provides a positive relative benefit when
selected objects are above the resolution floor and within-distribution
with the target scene. It provides no benefit when the failure mode
is domain shift, and it amplifies AP$_S$ degradation when selected
objects are sub-resolution. A practical diagnostic criterion follows:
run MinImageScorer, inspect the size and domain distribution of
high-score objects, and apply copy-paste augmentation only when
those objects satisfy both conditions above.

% ----------------------------------------------------------
\section{Conclusion}

We proposed MinImageScorer, a model-failure-driven difficulty scoring
framework validated on two agricultural detection benchmarks.
Controlled experiments show that score-guided copy-paste outperforms
random copy-paste on MinneApple ($+0.013$~AP) but both strategies
degrade below baseline, with the degradation confirmed by
feature-space analysis to originate from distributional shift
between augmented training images and the test set.
On GWHD, domain shift makes copy-paste ineffective under any
selection strategy.

The practical implication is that copy-paste augmentation should be
preceded by failure-mode diagnosis: it is effective for
underrepresentation (COCO-like datasets) but not for
resolution-limited or domain-shifted datasets. Future work includes
online augmentation to reduce distributional rigidity, and
domain-constrained pasting for multi-source benchmarks.

% ----------------------------------------------------------
\begin{thebibliography}{00}

\bibitem{hani2020minneapple}
N.~Hani, P.~Roy, and B.~Bhanu,
``MinneApple: A Benchmark Dataset for Apple Detection and
Segmentation,''
\textit{IEEE Robotics and Automation Letters},
vol.~5, no.~2, pp.~852--858, 2020.

\bibitem{david2021global}
E.~David \textit{et~al.},
``Global Wheat Head Detection (GWHD) Dataset,''
\textit{Plant Phenomics}, vol.~2021, 2021.

\bibitem{kisantal2019augmentation}
M.~Kisantal \textit{et~al.},
``Augmentation for Small Object Detection,''
in \textit{Proc. 9th Int. Conf. Advances in Computing,
Communication and Informatics (ICACCI)}, 2019.

\bibitem{ghiasi2021simple}
G.~Ghiasi \textit{et~al.},
``Simple Copy-Paste is a Strong Data Augmentation Method
for Instance Segmentation,''
in \textit{Proc. IEEE/CVF Conf. Computer Vision and Pattern
Recognition (CVPR)}, pp.~2918--2928, 2021.

% ============================================================
% [OPTIONAL] Additional references to consider:
% - R. Dufumier et al., GWHD challenge paper (if citing challenge)
% - T.-Y. Lin et al., Focal Loss / RetinaNet (ICCV 2017)
% - A. Shrivastava et al., OHEM (CVPR 2016)
% ============================================================

\end{thebibliography}

\end{document}
